\documentclass[10pt,a4paper]{amsart}
\usepackage[margin=19mm]{geometry}
\usepackage{amsmath,amssymb,mathtools}
\usepackage[scr=rsfs,cal=euler]{mathalfa}
\usepackage{xfrac}
\usepackage[unicode,hidelinks,bookmarks=false]{hyperref}
\newcommand{\ie}{i.e.\ }
\newcommand{\cf}{cf.\ }
\newcommand{\resp}{respectively\ }
\newcommand{\Subsection}{\subsection}
\providecommand{\nobreakdash}{\nobreak\hbox{-}}
\setlength{\emergencystretch}{2em}
\multlinegap0pt
\usepackage{bm}
\usepackage{bbm}
\usepackage{dsfont}
\usepackage{xspace}
\usepackage{cancel}
\usepackage{mathtools}
\usepackage{amsthm}
\makeatletter
\@ifundefined{theo}{\newtheorem{theo}{Theo}}{}
\@ifundefined{lemma}{\newtheorem{lemma}[theo]{Lemma}}{}
\makeatother
\makeatletter
\@ifundefined{proof}{\newenvironment{proof}{\noindent\textbf{Proof}.}{\hfill$\qed$}}{}
\newcommand{\sm}{\setminus}
\makeatother
\title[(Even hole, triangle)-free graphs revisited]{(Even hole, triangle)-free graphs revisited}
\author{Beatriz Martins, Nicolas Trotignon}
\date{Source: arXiv:2604.01816v1 (CC BY 4.0)}
\begin{document}
\maketitle

\noindent\emph{Source and licence.} (Even hole, triangle)-free graphs revisited. Version arXiv:2604.01816v1 (CC BY 4.0), distributed under the Creative Commons Attribution 4.0 International licence, as recorded in the arXiv licence metadata for this exact version and shown on the abstract page https://arxiv.org/abs/2604.01816v1. Full rights notice: licenses/arxiv-2604.01816-CC-BY-4.0.txt.\par\smallskip

\noindent\textit{Editorial scope.} Counted unit: the Lemma whose source label is \texttt{l:goodEquiv} in the article (source0.tex lines 911--920, source label \texttt{l:goodEquiv}), delivered together with the article's own complete original proof of it (source0.tex lines 922--966, one \texttt{proof} environment of 45 lines). No mathematics is written, repaired or re-derived: statement and proof are the article's own text line for line. Required context retained uncounted: none. Every definition, display and label of those lines is retained unchanged. Omitted from this package and not redistributed: 1--910, 921--921, 967--1794. Nothing inside a retained excerpt is omitted or altered: every retained body is delivered exactly as the article's lines are written, so no omission receipt of any kind accompanies this package. Citations: the retained excerpts carry no citation command at all, so no bibliography entry of the article is transcribed and no reference is redistributed. Reference adaptation: none was needed and none was made; every reference inside the delivered unit resolves inside it, so no unresolved reference or citation is produced. Printed slips kept literal: no printed word, symbol, hypothesis or number of the counted statement or of its proof was altered. Layout and class: the article's own class is \texttt{article} with packages including babel, authblk, amsmath, amssymb, amsthm, thmtools, graphicx, hyperref, tikz, comment, cleveref; the wrapper is the lane's pinned \texttt{amsart} editorial wrapper at 10pt on A4, which restates the theorem-like environments the retained text needs and adds the article's own macro definitions that the retained text needs (\texttt{\textbackslash sm}), with extra wrapper packages (bm, bbm, dsfont, xspace, cancel, mathtools). The delivered numbering is the unit's own. Assets: no retained span contains a figure environment, a graphics-inclusion command, a PDF-page inclusion, a table or a TikZ picture; no image file is copied into this package, the package declares no asset dependency and no image file is redistributed with it. Licence: version arXiv:2604.01816v1 of this article is distributed under the Creative Commons Attribution 4.0 International licence, as recorded in the arXiv licence metadata for this exact version and shown on the abstract page https://arxiv.org/abs/2604.01816v1. A full rights notice is delivered at licenses/arxiv-2604.01816-CC-BY-4.0.txt. Characters: every retained excerpt is pure ASCII, so the delivered engine, XeLaTeX with its default Latin Modern text font, reports no missing character.
\begin{lemma}
  \label{l:goodEquiv}
  Let $G$ be an atomic (theta, triangle)-free graph.
  A path $acb$ of~$G$ is a good $P_3$ if and only if it satisfies the following three conditions:    
  \begin{enumerate}
  \item\label{i:hole} some hole of $G$ goes through $a$, $c$ and $b$, 
  \item\label{i:center}  there is no wheel $(H, c)$ in $G$ such that $a, b \in H$ and
  \item\label{i:spoke} there is no wheel $(H, c')$ in $G$ such that $a, c, b \in H$ and $cc'\in E(G)$.
  \end{enumerate}
\end{lemma}

\begin{proof}
  Suppose first that $acb$ is good.  Since $G$ is atomic, there exists a path~$P$ from $a$ to $b$ in $G\sm c$, and this path has no internal vertex adjacent to $c$ since $acb$ is good. Hence, $acb$ and $P$ form a hole, so Condition~\ref{i:hole} is satisfied.  Conditions~\ref{i:center} and~\ref{i:spoke} are clearly satisfied for otherwise, there exists a path from $a$ to $b$ with internal vertices adjacent to $c$, a contradiction to $acb$ being good. 
  
  Let us prove the converse. So, suppose for a contradiction $acb$ is not good while Conditions~\ref{i:hole}, \ref{i:center} and~\ref{i:spoke} are satisfied. 
  So there exists an $ab$-path~$P$, not containing $c$ and containing at least one internal vertex adjacent to $c$.  We choose $P$ subject to the minimality of the  number $k\geq 1$ of internal vertices of $P$ adjacent to $c$.   
  
  Call \emph{sector} of $P$ any subpath of $P$ of length at least~1, whose ends are adjacent to $c$ and whose internal vertices are not (so $P$ has at least two sectors). 
  
  By Condition~\ref{i:hole}, some hole goes through $acb$. So there exists an $ab$-path~$Q$ in $G\sm c$ such that no internal vertex of $Q$ is adjacent to $c$. 
  
  Not all vertices of $Q$ are in $P$ because the interior of $Q$ contains no neighbour of $c$. 
  Since $a$ and $b$ are in distinct sectors of $P$, there exists a subpath $Q'= a'\dots b'$ of $Q$ such that $a'$ and $b'$ are vertices of distinct sectors $S'$ and $T'$ of $P$, and the interior of $Q'$ is disjoint from $P$. Note that $a'$ (resp.\ $b'$) is either one of $a$ or $b$, or is an internal vertex of $S'$ (resp.\ $T'$). 

  Consider a subpath $R= x \dots y$ of $Q'$ and two distinct sectors $S = s\dots r$ and $T = r'\dots t$ of $P$ such that $x$ has neighbours in $S\sm T$, $y$ has neighbours in $T\sm S$ and no sector of $P$ contains $N_P(x) \cup N_P(y)$. 
  The existence of $R$ comes from the fact that the interior of $Q'$  satisfies these conditions.  
  Suppose that $S$, $T$ and $R$ are chosen subject to the minimality of $R$ (possibly $x=y$). 
  
  Up to symmetry, we suppose that $a$, $s$, $r$, $r'$, $t$ and $b$ appear in this order along $P$.  
  Let $x'$ (resp.\ $x''$) be the neighbour of $x$ in $S$ closest to $s$ (resp.\ to $r$) along $S$.  Let $y'$ (resp.\ $y''$) be the neighbour of $y$ in $T$ closest to $r'$ (resp.\ to $t$) along $T$. 
  
  Since $R$ is a subpath of the interior of $Q'$, it contains neither $c$ nor neighbours of $c$ and $R$ is disjoint from $P$.  Also, by the minimality of $R$, an internal vertex of $R$ has no neighbours in $P$, except possibly $r$ if $r=r'$.  Finally, no internal vertex of $r P r'$ has a neighbour in $R$,  except possibly if $x=y$.  
 
  Consider the path $Z = a P x' x R y y'' P b$.  
  We claim that some neighbour of~$c$  in $rPr'$ is not in $Z$.  
  Otherwise $r=r'$ or $rPr'$ is a sector of $P$, and $N_P(x) \cup N_P(y)$ is included in some sector of $P$, a contradiction to the definition of~$R$. 
  Hence, $Z$  contains less neighbours of $c$ than~$P$.  So, the interior of $Z$ contains no neighbour of $c$ for otherwise, $Z$ contradicts the minimality of $k$.  
  This implies that $S$ is the sector of $P$ that contains $a$ (so $s=a$), $T$ is the sector of $P$ that contains $b$ (so $t=b$), $x' \neq  r$ and $y'' \neq  r'$. 

  Suppose $x'\neq x''$.  
  Then $x x' P a c$, $x x'' P r c$ and $x R y y'' P b c$ form an $xc$-theta unless $r'=r$, $y'=y''$ and $y'r \in E(G)$.  
  But then, $y' r c$, $y' P b c$ and $y'y R x x' P a c$ form a $y'c$-theta unless $x''=r$. 
  In this last case, $(aPx'xRyy'Pbca, r)$ is a wheel, a contradiction to Condition~\ref{i:spoke} (note that in this case, $P$, $R$ and $c$ form a daisy with two petals). 
  This proves that $x'=x''$, and symmetrically $y'=y''$ (recall that $x' \neq  r$ and $y'' \neq  r'$).
  
  Suppose that $x'\neq a$.  
  Then, the three paths $x' P a c$, $x' P r c$ and $x' x R y y' P b c$ form a theta unless $r = r'$ and $y'r \in E(G)$.  In this last case, $y'\neq b$, so a symmetric argument proves that $x'r\in E(G)$.  
  Then, the hole $a P x' x R y y' P b c a$ is the rim of wheel centred at $r$, a contradiction to Condition~\ref{i:spoke}.  Hence, $x'=a$, and symmetrically, we can prove $y'=b$. 
  
  Suppose that $r=r'$.  If $r$ has some neighbours in $R$, then $rca$, $rPa$ and a shortest path from $r$ to $a$ in $\{r, a\} \cup R$ form a theta, a contradiction. So, $r$ has no neighbours in $R$.  
  Hence, $axRybPa$ is a hole, and therefore the rim of a wheel centred at $c$.  
  This is a contradiction to Condition~\ref{i:center}.  So, $r\neq r'$.  By the minimality of $R$, it follows that no internal vertex of $R$ has a neighbour in $P$.   
  
  Suppose that some vertex of $rPr'$  has some neighbour in $R$.  By the minimality of $R$, this implies that $R$ has length~0, so $x=y$.  Then, $xac$, $ybc$ and some shortest path from $x$ to $c$ in $\{x, c\} \cup rPr'$ form an $xc$-theta, a contradiction. Hence, no vertex of $rPr'$ has some neighbour in $R$. 
  Hence, $axRybPa$ is a hole,  and therefore the rim of a wheel centred at $c$.  This is a contradiction to Condition~\ref{i:center}. 
\end{proof}

\end{document}
